% !TeX spellcheck = en_GB
% !TeX program = lualatex
%
% v 2.3  Feb 2019   Volker RW Schaa
%		# changes in the collaboration therefore updated file "jacow-collaboration.tex"
%		# all References with DOIs have their period/full stop before the DOI (after pp. or year)
%		# in the author/affiliation block all ZIP codes in square brackets removed as it was not %         understood as optional parameter and ZIP codes had bin put in brackets
%       # References to the current IPAC are changed to "IPAC'19, Melbourne, Australia"
%       # font for ‘url’ style changed to ‘newtxtt’ as it is easier to distinguish "O" and "0"
%
\documentclass[a4paper,
               %boxit,        % check whether paper is inside correct margins
               %titlepage,    % separate title page
               %refpage       % separate references
               biblatex,     % biblatex is used
               %keeplastbox,   % flushend option: not to un-indent last line in References
               %nospread,     % flushend option: do not fill with whitespace to balance columns
               %hyphens,      % allow \url to hyphenate at "-" (hyphens)
               %xetex,        % use XeLaTeX to process the file
               %luatex,       % use LuaLaTeX to process the file
               ]{jacow}
%
% ONLY FOR \footnote in table/tabular
%
\usepackage{pdfpages,multirow,ragged2e} %
%
% CHANGE SEQUENCE OF GRAPHICS EXTENSION TO BE EMBEDDED
% ----------------------------------------------------
% test for XeTeX where the sequence is by default eps-> pdf, jpg, png, pdf, ...
%    and the JACoW template provides JACpic2v3.eps and JACpic2v3.jpg which
%    might generates errors, therefore PNG and JPG first
%
\makeatletter%
	\ifboolexpr{bool{xetex}}
	 {\renewcommand{\Gin@extensions}{.pdf,%
	                    .png,.jpg,.bmp,.pict,.tif,.psd,.mac,.sga,.tga,.gif,%
	                    .eps,.ps,%
	                    }}{}
\makeatother

% CHECK FOR XeTeX/LuaTeX BEFORE DEFINING AN INPUT ENCODING
% --------------------------------------------------------
%   utf8  is default for XeTeX/LuaTeX
%   utf8  in LaTeX only realises a small portion of codes
%
\ifboolexpr{bool{xetex} or bool{luatex}} % test for XeTeX/LuaTeX
 {}                                      % input encoding is utf8 by default
 {\usepackage[utf8]{inputenc}}           % switch to utf8

\usepackage[USenglish]{babel}
\usepackage{stackengine}
%
% if BibLaTeX is used
%
\ifboolexpr{bool{jacowbiblatex}}%
 {%
  \addbibresource{MOP019.bib}
%  \addbibresource{biblatex-examples.bib}
 }{}
%\listfiles

%%
%%   Lengths for the spaces in the title
%%   \setlength\titleblockstartskip{..}  %before title, default 3pt
%%   \setlength\titleblockmiddleskip{..} %between title + author, default 1em
%%   \setlength\titleblockendskip{..}    %afterauthor, default 1em

\begin{document}

\title{BUNCH DURATION MEASUREMENTS IN THE APS-U BOOSTER\thanks{Work supported by the U.S. D.O.E.,Office of Science, Office of Basic Energy Sciences, under contract number DE-AC02- 06CH11357.}}

\author{J. C. Dooling\thanks{dooling@anl.gov}, J. R. Calvey, O. Mohsen, Argonne National Laboratory, Lemont, USA \\
}

%A. Brill, K. Wootton
	
\maketitle

%
\begin{abstract}

 We present the results of time-based, bunch length measurements in the Advanced Photon Source Upgrade booster synchrotron using the bunch duration monitor (BDM) optical diagnostic.  The BDM diagnostic is based on the detection of visible-wavelength synchrotron radiation.  The detector is a metal-semiconductor-metal device followed by 42 dB of wide-band amplifier gain.  Bunch duration is determined by de-convolving the raw output signal with the circuit’s impulse response function.  The BDM allows measurement of bunch duration over virtually the entire booster ramp.  Defocusing in the optical path was necessary to overcome thermal steering from the in-tunnel mirror.  Also, the effects of detector saturation must be considered to ensure a linear response. Presently, the booster increases bunch energy from 425 MeV to 6 GeV.  Booster charge varies from 5 nC to 15 nC depending on storage ring operating modes.  BDM data reveal that the bunch undergoes large longitudinal oscillations shortly after injection into the booster.  The longitudinal oscillations are compared with elegant simulations.  These oscillations are a source of injection loss especially at higher charge.
 
\end{abstract}
%\end{document}
%\section{Comments}
%\begin{itemize}
%    \item Signal to noise ratio is small at injection when bunch is spread out (duration large) and energy is low (gamma low)
%    \item Bunch duration from the PAR $\geq600$\,ps at high charge
%    \item Synchrotron rad intensity as a function of energy (gamma)
%\end{itemize}

\section{Introduction}

The Advanced Photon Source Upgrade (APS-U) booster synchrotron provides the major boost in energy from the linac at 400-450 MeV to storage ring (SR) at 6 GeV. 
%Timing of the injected pulse into the SR is critical for swap-out injection.
Bunch duration is an important parameter for efficient booster operation, especially at injection~\cite{calvey:ipac21-mopab045}.  

\section{Experimental Arrangement}

The booster BDM is located within the Synchrotron Light Monitor 2 (SLM2) enclosure in the A015 optics lab directly above the booster tunnel.  The BDM front-end electronics is similar to that described earlier for the APS particle accumulator ring~\cite{dooling:ipac18-wepaf006}.
An in-tunnel mirror reflects synchrotron radiation up out of the booster tunnel to an optics station above the shield wall.

\subsection{Saturation Effects}

Broadening of the pulse duration measurement is observed when the signal is strong, leading to incorrect BDM results~\cite{dooling:napac22-tupa22}.  Saturation effects can be mitigated if de-convolved pulse amplitudes are kept below 8 V.  We accomplished this requirement by defocusing the image on the detector as discussed in the next subsection.

\subsection{Defocusing Optics}

The active region size of the MSM detector is 0.2 mm by 0.2 mm.  When the detector was placed near the focus of the final focusing lens, in addition to saturation, motion of the image spot due to thermal distortion of the in-tunnel mirror made measurements challenging.  Initially, an ND1.0 neutral density filter was placed between the lens and the detector.  In order to minimize the thermal effects, the filter was removed and the detector was placed closer to the lens.  The spot on the detector was then defocused and larger; this significantly improved availability of the diagnostic by reducing the need for realignment.  The modification is shown schematically in Fig.~\ref{fig:booBDMoptics}.
\begin{figure}
    \centering
    \includegraphics[width=0.8\linewidth]{figures/MOP019_f1.png}
    \caption{Optical layout of the booster SLM2 enclosure with the defocus modification of the BDM (D/A) line indicated.}
    \label{fig:booBDMoptics}
\end{figure}

\section{Measurements}

The BDM diagnostic with modified optics was employed to measure bunch duration throughout the booster energy ramp at different charge levels.  With the modified optics, the diagnostic can routinely gather data during operations. Signal optimization using a 2-axis picomotor-controlled mirror mount is done only occasionally now; for example, when changing from low- to high-charge operations or significantly changing the repetition rate.

Due to the inverse-square dependence of synchrotron radiation power on bending radius, $\rho$~\cite{Walker:398429}, the synchrotron radiation power from the booster at injection is much less than that from the particle accumulator ring (PAR), $\rho_{\text{PAR}}=1.02\,$m versus $\rho_{\text{Boo}}=33.3\,$m; thus the booster BDM signal at injection tends to be weaker than that from the PAR BDM at PAR extraction.  

The PAR collects charge from the s-band linac as a series of 30-Hz pulses into a single bunch ($h=1$, fundamental rf frequency, $f_0=9.78$\,MHz), then compresses the bunch with a 12$^{th}$-harmonic cavity for injection to the booster at a maximum rate of 1 Hz.  In all cases presented here, linac energy and hence PAR and booster injection energy, were 425 MeV. Depending on stored charge, the rms bunch duration at extraction from the PAR ranges from 300-600 ps; a recent measurement is presented in Fig.~\ref{fig:PARextr}.
\begin{figure}[t!]
    \centering
    \includegraphics[width=0.8\linewidth]{figures/MOP019_f2.png}
    \caption{Measured PAR rms bunch duration shortly before extraction versus PAR-to-Booster (PTB) charge.}
    \label{fig:PARextr}
\end{figure}
PAR BDM measurements have been shown to agree with simultaneous streak camera data~\cite{dooling:ipac18-wepaf006}.

RMS bunch duration from BDM data (both PAR and Booster) is determined by deconvolution of the raw data waveform with the impulse response function for the detector and amplifier, the resulting pulse shape is then fit with a Gaussian function.

\subsection{Injection Oscillations}

At higher charge, signal-to-noise (S/N) is sufficient to capture and measure the bunch during the first few turns after booster injection; an example is shown in Fig.~\ref{fig:booInj}.  A significant increase in signal amplitude is observed after injection.  The nominal booster revolution period is $1.2275\,\mu$s ($h=432$, single bunch). Small modifications are made to the 352 MHz drive frequency to properly time the bunch injected into the storage-ring for swap-out operations~\cite{calvey:ipac24-tupg07}.
%\begin{figure}
%    \centering
%    \includegraphics[width=0.85\linewidth]{figures/BDM_080_005_7.proc_MedResTrim.png}
%    \caption{Booster BDM waveform showing the first few turns after injection. In this case, injection occurs between 3 and 4 $\mu$s.  PTB charge was 14.1 nC.}
%    \label{fig:booInj}
%\end{figure}

\begin{figure}[h!]
    \centering
    \def\big{\includegraphics[width=0.85\linewidth]{figures/MOP019_f3a.png}}
    \def\little{\includegraphics[width=0.45\linewidth]{figures/MOP019_f3b.png}}
    \stackinset{r}{54pt}{b}{93pt}{\little}{\big}
    \caption{Booster BDM waveform showing the first few turns after injection.  PTB charge was 14.1 nC.  The inset shows charge during the full booster cycle with a PTB input of 11.7 nC. Charge drops to 10.4 nC shortly after injection indicating a loss of 1.3 nC.} 
    %In this case, injection occurs between 3 and 4 $\mu$s.  
    \label{fig:booInj}
\end{figure}

Figure~\ref{fig:booInjOsc} shows a longer BDM waveform recorded on an oscilloscope with deeper memory; here variation in peak amplitude is seen to oscillate for several dozen turns (50\,$\mu$s). PTB charge in this case was 15.3 nC.

\begin{figure}[t]
    \centering
    \includegraphics[width=0.8\linewidth]{figures/MOP019_f4.png}
    \caption{Booster BDM waveform from a longer record showing bunch length oscillation shortly after injection.  In this case, PTB charge was 15.3 nC.}
    \label{fig:booInjOsc}
\end{figure}

An example of the deconvolved waveform at injection is given in Fig.~\ref{fig:deconvTurn1}; this corresponds to the first pulse seen in Fig.~\ref{fig:booInj}.  
\begin{figure}[b!]
    \centering
    \includegraphics[width=0.8\linewidth]{figures/MOP019_f5.png}
    \caption{Deconvolved bunch waveform and Gaussian fit to the first turn in the booster.}
    \label{fig:deconvTurn1}
\end{figure}
Here $\sigma_t=315$\,ps from a Gaussian fit and is reasonably consistent with a FWHM calculation of 708 ps ($2\sqrt{2\ln{2}}\sigma_t=741$\,ps).  Though the S/N is relatively low, similar durations are measured in other cases where the injected bunch is observed, see Table~\ref{tab:InjDur}.  
\begin{table}[b]
    \centering
    \caption{Injection charge and Gaussian fit and FWHM bunch duration for first and second turns in the booster.}
    \begin{tabular}{|c|c|c|c|c|c|}
    \hline
    \multirow{1}{*}{delay} & PTB &   
%    \cline{2-7}
    \multicolumn{2}{ c |}{ Turn 1} & \multicolumn{2}{ c |}{Turn 2} \\
 ($\mu$s)   & charge & $\sigma_t$ & FWHM  & $\sigma_t$ & FWHM  \\
  \_case      & (nC) & (ps) & (ps)   & (ps) & (ps)  \\
                \hline
 10\_2 & 13.1   &  352 & 752 & 328  &  855 \\
 05\_3 & 13.1  &  316  & 726 & 289  & 708 \\
 05\_7 & 14.0  & 315   & 708 & 332  &  823 \\
 \hline
    \end{tabular}
    \label{tab:InjDur}
\end{table}
Also shown in the table are durations measured for the second turn which show little change from the first. The average of all rms values in the table yields, $\sigma_t=322\pm22$\,ps.
On the other hand, where peak signal is high and S/N is good, the bunch duration drops dramatically as presented in Fig.~\ref{fig:peakSig}.
\begin{figure}[h]
    \centering
    \includegraphics[width=0.8\linewidth]{figures/MOP019_f6.png}
    \caption{Deconvolved bunch waveform and Gaussian fit shortly after injection near peak signal strength.}
    \label{fig:peakSig}
\end{figure}

The bunch duration in terms of both $\sigma_t$ and FWHM for the waveform presented in Fig.~\ref{fig:booInjOsc} is given in Fig.~\ref{fig:durOscill}.  Note that the FWHM algorithm can be fooled by data spikes which took place in the data set shortly after 30\,$\mu$s.  Fig.~\ref{fig:durOscill} clearly shows large bunch duration oscillations shortly after injection into the the booster.  This longitudinal mismatch likely leads to particle loss at injection.
\begin{figure}
    \centering
    \includegraphics[width=0.88\linewidth]{figures/MOP019_f7.png}
    \caption{$\sigma_t$ and FWHM for the waveform shown in Fig.~\ref{fig:booInjOsc}}
    \label{fig:durOscill}
\end{figure}

\subsection{Comparison with simulations}

These results can be compared with particle tracking simulations.  The simulations include a model of the booster impedance, beam loading in the rf cavities, and incoming beam parameters derived from measurements at PAR extraction~\cite{calvey:ipac21-mopab045}.  The longitudinal profile of the beam is output at every turn, and fitted with a Gaussian function.  This can be directly compared with the BDM measurements (Fig.~\ref{fig:simComp}).  The simulation also shows a large variation in bunch length, which is a result of the mismatched bunch from the PAR tumbling inside the booster bucket.  This tumbling is a contributor to losses at booster injection~\cite{calvey:napac25-MOYN01}.  The measured oscillation period was reduced by 11\% to match the simulation result.

\begin{figure}
    \centering
    \includegraphics[width=0.9\linewidth]{figures/MOP019_f8.png}
    \caption{Comparison of measured (Fig.~\ref{fig:booInjOsc}) and simulated bunch duration shortly after booster injection, 15~nC.}
    \label{fig:simComp}
\end{figure}

\subsection{Booster Ramp Scan}

BDM measurements were carried out over the full booster ramp cycle ($\approx190$\,ms) in 10 ms steps after injection.  Comparison of moderate and high-charge scans are given in Fig.~\ref{fig:booScans}.  Charge is recorded approximately mid-way during the ramp.

\begin{figure}[h]
    \centering
    \def\big{\includegraphics[width=0.95\linewidth]{figures/MOP019_f9a.png}}
    \def\little{\includegraphics[width=0.45\linewidth]{figures/MOP019_f9b.png}}
    \stackinset{r}{35pt}{b}{63pt}{\little}{\big}
    \caption{BDM measurements in the booster during the energy ramp from 425 MeV to 6 GeV at both medium and high-charge operation.  The inset zooms in on the shorter bunch durations measured at later times.  In this plot, the initial 80 ms delay has been subtracted from the time axis. Extraction occurs near 190 ms; data at 200 ms are from cycles not extracted.}
    \label{fig:booScans}
\end{figure}

\section{Summary}

Defocusing of the booster BDM optics has reduced sensitivity to thermal steering from the in-tunnel mirror and facilitated on-line measurement of bunch length throughout the acceleration cycle.  Strong longitudinal oscillations at injection are observed and largely agree with simulations.  This longitudinal mismatch leads to significant beam loss at injection.

\section{ACKNOWLEDGMENTS}

We thank R. Soliday and H. Shang for assistance with analysis and data acquisition scripts. We also are grateful to W. Cheng for providing resources and hardware.
%Any acknowledgement should be in a separate section directly preceding
%the \textbf{REFERENCES} or \textbf{APPENDIX} section.



%
% only for "biblatex"
%
\ifboolexpr{bool{jacowbiblatex}}%
	{\printbibliography}%
	%{%
	% "biblatex" is not used, go the "manual" way
	
	%\begin{thebibliography}{99}   % Use for  10-99  references
%	\begin{thebibliography}{9} % Use for 1-9 references
	
%	\bibitem{jacow-help}
%		JACoW,
%		\url{http://www.jacow.org}
	
%	\bibitem{IEEE}
%		\textit{IEEE Editorial Style Manual},
%		IEEE Periodicals, Piscataway,
%		NJ, USA, Oct. 2014, pp. 34--52.

%	\bibitem{journal-abbreviations}
%	\url{https://woodward.library.ubc.ca/researchhelp/journal-abbreviations/}

%	\end{thebibliography}
%}  end \ifboolexpr
%
% for use as JACoW template the inclusion of the ANNEX parts have been commented out
% to generate the complete documentation please remove the "%" of the next two commands
% 
%%%\newpage

%%%\include{annexes-A4}
\clearpage
\centering
\fbox{\begin{minipage}{35em}
{\footnotesize The submitted manuscript has been created by UChicago Argonne, LLC, Operator of
Argonne National Laboratory (“Argonne”). Argonne, a U.S. Department of Energy
Office of Science laboratory, is operated under Contract No. DE-AC02-06CH11357.
The U.S. Government retains for itself, and others acting on its behalf, a paid-up
nonexclusive, irrevocable worldwide license in said article to reproduce, prepare
derivative works, distribute copies to the public, and perform publicly and display
publicly, by or on behalf of the Government. The Department of Energy will provide
public access to these results of federally sponsored research in accordance with the
DOE Public Access Plan. http://energy.gov/downloads/doe-public-access-plan}
\end{minipage}
}

\end{document}



